\documentclass[12pt]{article}

\usepackage[utf8]{inputenc}
\usepackage{amsmath, amssymb, amsthm}
\usepackage{geometry}
\geometry{a4paper, margin=2.5cm}

\title{Structura Submodulelor Libere peste un PID}
\author{}
\date{}

\theoremstyle{definition}
\newtheorem{definition}{Definiție}

\theoremstyle{plain}
\newtheorem{theorem}{Teoremă}

\begin{document}

\maketitle

\begin{theorem}[Structura submodulelor libere peste un PID]
Fie $R$ un domeniu principal (PID), $F$ un $R$-modul liber de rang $n$ și
$E \subseteq F$ un submodul de rang $m$. Atunci există o bază


\[
u_1, \dots, u_n
\]


a lui $F$ și elemente nenule


\[
d_1, \dots, d_m \in R
\]


cu proprietățile:
\begin{enumerate}
    \item[(i)] $d_1 \mid d_2 \mid \cdots \mid d_m$;
    \item[(ii)] elementele
    

\[
    d_1 u_1,\; d_2 u_2,\; \dots,\; d_m u_m
    \]


    formează o bază a submodulului $E$.
\end{enumerate}
\end{theorem}

\begin{theorem}[Formă normală Smith pentru submodule]
În ipotezele teoremei precedente, există matrici inversabile
$U \in GL_n(R)$ și $V \in GL_m(R)$ astfel încât matricea de includere
$E \hookrightarrow F$ se poate scrie sub forma diagonală


\[
\operatorname{diag}(d_1, \dots, d_m, 0, \dots, 0),
\]


unde $d_1 \mid d_2 \mid \cdots \mid d_m$.
\end{theorem}

\end{document}
